% Leaf Target Sufficiency Lemma

\subsection{Leaf target sufficiency}\label{sec:leaf-lemma}

The solver in Section~\ref{sec:computational} tests only the leaf targets $u_L$ and~$u_R$.

\begin{lemma}[Leaf target sufficiency]\label{lem:leaf}
Let $G=T(3,s_1,s_2)$ with leaves $u_L$ and $u_R$. Then
\[
\pi(G)=\max\bigl\{\pi(G,u_L),\,\pi(G,u_R)\bigr\}.
\]
Equivalently, if every configuration with $S$ pebbles can pebble both leaves,
then every such configuration can pebble every vertex of~$G$.
\end{lemma}

\begin{lemma}[Tail propagation]\label{lem:tail-prop}
Let $P=(w_0,w_1,\ldots,w_k)$ be a path in $G$ with $w_k$ a leaf of~$G$.
For every $i\in\{0,1,\ldots,k\}$,
\[
\pi(G,w_i)\le \pi(G,w_k).
\]
In particular, with $P_L=(a,\ldots,u_L)$ and $P_R=(b,\ldots,u_R)$ one has
$\pi(G,v)\le \pi(G,u_L)$ for every vertex $v$ on the left tail and
$\pi(G,v)\le \pi(G,u_R)$ for every vertex $v$ on the right tail.
\end{lemma}

\begin{lemma}[First double on a path]\label{lem:first-double}
Let $P=(w_0,w_1,\ldots,w_k)$ be a path and let $D$ be a configuration that
pebbles~$w_i$ for some $i\in\{0,1,\ldots,k\}$.
If some reachable state satisfies $p(w_i)\ge 2$, then $D$ pebbles~$w_i$.
Otherwise $p(w_i)\le 1$ in every reachable state, and the first move that
places a pebble on~$w_i$ is either $w_{i-1}\to w_i$ (when $i\ge 1$) or
$w_{i+1}\to w_i$ (when $i\le k-1$); in either case the prefix ending
immediately before that move already places at least two pebbles on the source
vertex, so $D$ pebbles that source as well.
\end{lemma}

\begin{proof}
If a reachable state has $p(w_i)\ge 2$, then either $w_i\to w_{i-1}$ or
$w_i\to w_{i+1}$ is legal toward pebbling a neighbor, and the claim is
immediate from the definition of ``$D$ pebbles~$w_i$''.

Assume $p(w_i)\le 1$ in every reachable state.
Let $m$ be the first move in a pebbling sequence for~$w_i$ that places a pebble
on~$w_i$.
Because $p(w_i)=0$ before~$m$, the move $m$ is either $w_{i-1}\to w_i$ or
$w_{i+1}\to w_i$.
Each requires at least two pebbles on its source immediately beforehand, and
the prefix ending before~$m$ pebbles that source.
\end{proof}

\begin{lemma}[Left junction supply for the apex]\label{lem:left-junction}
Let $G=T(3,s_1,s_2)$ and let $w_1$ denote the neighbor of~$a$ on the left
tail ($w_1=u_L$ when $s_1=1$).
Fix a configuration~$p$ that can pebble~$a$.
If every reachable configuration has $p(c)=0$ and $p(b)\le 1$, then some
reachable configuration~$p'$ satisfies $p'(a)\ge 2$.
\end{lemma}

\begin{proof}
If some reachable state already has $p(a)\ge 2$, we are done.
Assume $p(a)\le 1$ in every reachable state.
Because $p$ pebbles~$a$, some sequence ends with a pebble on~$a$.

The first move~$m$ that places a pebble on~$a$ cannot be $c\to a$ or $b\to a$,
because $p(c)=0$ and $p(b)\le 1$ in every reachable state.
Hence $m=w_1\to a$, and the prefix ending immediately before~$m$ already has
$p(w_1)\ge 2$ by Lemma~\ref{lem:first-double}.

If $p(a)=0$ before~$m$, then after~$m$ we have $p(a)=1$.
Because $p(a)\le 1$ always, a second move must place another pebble on~$a$;
it is again $w_1\to a$, and immediately beforehand $p(w_1)\ge 2$.
Performing this second move yields a reachable state with $p(a)\ge 2$.
\end{proof}

\begin{proof}[Proof of Lemma~\ref{lem:tail-prop}]
Fix $S=\pi(G,w_k)$ and let $D$ be any configuration with $|D|=S$.
We prove by downward induction on $i$ that $D$ pebbles~$w_i$.

The base $i=k$ is immediate.
For $i<k$, assume $D$ pebbles~$w_{i+1}$.
If $D$ already pebbles~$w_i$, we are done.
Otherwise $p(w_i)=0$.

\smallskip\noindent
\textbf{Cases 1--2.}
If $p(w_{i+1})=0$, the first move placing a pebble on~$w_{i+1}$ is $w_i\to w_{i+1}$,
and its prefix pebbles~$w_i$.
If $p(w_{i+1})\ge 2$, then $w_{i+1}\to w_i$ pebbles~$w_i$.

\smallskip\noindent
\textbf{Case 3: $p(w_{i+1})=1$ and $p(w_i)=0$.}
Suppose $D$ does not pebble~$w_i$.

If some sequence raises $p(w_{i+1})$ to~$2$, let $m$ be the first such move.
Then $m$ is $w_i\to w_{i+1}$ or $w_{i+2}\to w_{i+1}$ (when $i+2\le k$).
In the first case the prefix before~$m$ already places two pebbles on~$w_i$.
In the second case, apply the same first-double chase on the suffix
$(w_{i+1},\ldots,w_k)$: when $i=k-1$ the only possible first accumulation is
$w_{k-1}\to w_k$, whose prefix pebbles $w_{k-1}=w_i$; for $i<k-1$, iterate until
a prefix pebbles~$w_i$.
Both contradict the assumption that $D$ does not pebble~$w_i$.

Hence $p(w_{i+1})\le 1$ in every reachable state, so neither
$w_{i+1}\to w_i$ nor $w_{i+1}\to w_{i+2}$ is ever legal.

If $i=k-1$, then $w_{i+1}=w_k$ is the leaf.
The lone pebble on~$w_k$ cannot be inert, for otherwise $|D|-1$ pebbles would
suffice to pebble~$w_k$, contradicting $S=\pi(G,w_k)$.
Some sequence therefore executes $w_{k-1}\to w_k$, and the prefix before that move
places at least two pebbles on~$w_{k-1}=w_i$ (Lemma~\ref{lem:first-double}).

If $i<k-1$, no sequence can use $w_{k-1}\to w_k$, for that would force
$p(w_{i+1})\ge 2^{\,k-i-1}$ by Lemma~\ref{lem:path-lb} on the subpath
$(w_{i+1},\ldots,w_k)$, contradicting the paragraph above.
Because $D$ pebbles~$w_{k-1}$, the first move placing a pebble on~$w_{k-1}$ is either
$w_{k-2}\to w_{k-1}$ or $w_k\to w_{k-1}$; in either case the prefix before that move
pebbles a neighbor closer to~$w_i$, and repeating the first-double chase along
$(w_i,\ldots,w_{k-1})$ eventually yields a prefix that pebbles~$w_i$ or a move
$w_i\to w_{i+1}$ whose prefix pebbles~$w_i$.

In every subcase, $D$ pebbles~$w_i$.
Thus $\pi(G,w_i)\le S=\pi(G,w_k)$ for every $i\in\{0,\ldots,k\}$.
\end{proof}

\begin{lemma}[Triangle access to the apex]\label{lem:triangle-apex}
Let $G=T(3,s_1,s_2)$ and let $D$ be a configuration with $|D|=\pi(G,u_R)$
that can pebble~$u_R$.
If $p(c)=0$ in every reachable state and $p(b)\le 1$ in every reachable state,
then $D$ can pebble~$c$.
\end{lemma}

\begin{proof}
If some reachable state has $p(b)\ge 2$, then $b\to c$ places a pebble on~$c$.
Assume $p(b)\le 1$ in every reachable state, so $b\to c$ is never legal and any
pebbling of~$c$ must use $a\to c$, which requires a reachable state with
$p(a)\ge 2$.
Because $\pi(G,a)\le \pi(G,u_R)=|D|$, configuration~$D$ can pebble~$a$.
Since $p(c)=0$ and $p(b)\le 1$ in every reachable state,
Lemma~\ref{lem:left-junction} gives a reachable state~$\sigma$ with
$\sigma(a)\ge 2$, and then $a\to c$ pebbles~$c$.

In every subcase, $D$ can pebble~$c$.
\end{proof}

\begin{lemma}[Apex bound]\label{lem:apex-bound}
The triangle apex satisfies
\[
\pi(G,c)\le \max\bigl\{\pi(G,u_L),\,\pi(G,u_R)\bigr\}.
\]
\end{lemma}

\begin{proof}
By symmetry, assume $\pi(G,u_R)\ge \pi(G,u_L)$.
Let $D$ be any configuration with $|D|=\pi(G,u_R)$ pebbles.
Then $D$ can pebble~$u_R$.

If $D$ already places a pebble on~$c$, then $\pi(G,c)\le \pi(G,u_R)$.

Otherwise $p(c)=0$.
We show that $D$ can pebble~$c$.

\smallskip\noindent
\textbf{Case 1: $p(u_R)=0$ in~$D$.}
Choose a pebbling sequence from~$D$ that places a pebble on~$u_R$.
Since $p(u_R)=0$, Lemma~\ref{lem:path-lb} applied to the right tail yields a reachable
configuration~$\sigma$ with $\sigma(b)\ge 2^{s_2}$.
If $\sigma(c)\ge 1$, then $D$ pebbles~$c$.
Otherwise the move $b\to c$ is legal in~$\sigma$ and places a pebble on~$c$.

\smallskip\noindent
\textbf{Case 2: $p(u_R)\ge 1$ in~$D$.}
Then $D$ already pebbles~$u_R$ by the empty sequence.
If some reachable state~$\sigma$ satisfies $\sigma(b)\ge 2$ or $\sigma(a)\ge 2$,
then a single move $b\to c$ or $a\to c$ pebbles~$c$.
Otherwise $p(a)\le 1$, $p(b)\le 1$, and $p(c)=0$ in every reachable state.
These are exactly the hypotheses of Lemma~\ref{lem:triangle-apex}, so $D$ can
pebble~$c$ by that lemma.

\smallskip\noindent
In both cases, every configuration with $\pi(G,u_R)$ pebbles can pebble~$c$, so
$\pi(G,c)\le \pi(G,u_R)=\max\{\pi(G,u_L),\pi(G,u_R)\}$.
\end{proof}

\begin{proof}[Proof of Lemma~\ref{lem:leaf}]
By Lemma~\ref{lem:tail-prop}, every vertex on the left tail, including~$a$, satisfies
$\pi(G,v)\le \pi(G,u_L)$; every vertex on the right tail, including~$b$, satisfies
$\pi(G,v)\le \pi(G,u_R)$.
By Lemma~\ref{lem:apex-bound},
$\pi(G,c)\le \max\{\pi(G,u_L),\pi(G,u_R)\}$.
Therefore every vertex $t$ satisfies
\[
\pi(G,t)\le \max\bigl\{\pi(G,u_L),\,\pi(G,u_R)\bigr\}.
\]
The reverse inequality is immediate because $u_L$ and $u_R$ are vertices of~$G$.
\end{proof}

\begin{corollary}[Solver pruning]\label{cor:pruning}
The pebbling number of $T(3,s_1,s_2)$ equals the minimum $S$ such that every
configuration with $S$ pebbles can pebble both leaves.
\end{corollary}

\begin{table}[ht]
\centering
\caption{Per-target pebbling numbers from exhaustive search
(\texttt{verify\_leaf\_lemma.py}).}
\label{tab:leaf-verify}
\begin{tabular}{lrrrrrrr}
\hline
Graph & $d$ & $\pi(G,c)$ & $\pi(G,a)$ & $\pi(G,b)$ & $\pi(G,u_L)$ & $\pi(G,u_R)$ & $\pi(G)$ \\
\hline
$T(3,1,1)$ & 3 & 7  & 6  & 6  & 9  & 9  & 9  \\
$T(3,1,2)$ & 4 & 11 & 10 & 8  & 17 & 17 & 17 \\
$T(3,2,2)$ & 5 & 15 & 12 & 12 & 33 & 33 & 33 \\
$T(3,1,3)$ & 5 & 19 & 18 & 12 & 33 & 33 & 33 \\
\hline
\end{tabular}
\end{table}

Each row agrees with Lemmas~\ref{lem:leaf} and~\ref{lem:apex-bound}.
On $T(3,1,1)$, $\pi(G,c)=7$ lies strictly below both leaf values.
Output: \path{data/leaf_lemma_verification.json}.
